\documentclass[a4paper,11pt]{article}
\usepackage[utf8]{inputenc}
\usepackage[T1]{fontenc}
\usepackage[english]{babel}
\usepackage[left=2.5cm,right=2.5cm,top=2cm,bottom=2cm]{geometry}
\usepackage{hyperref}
\usepackage{xcolor}
\usepackage{microtype}
\usepackage{lmodern}

\definecolor{primary}{RGB}{0, 45, 114}
\hypersetup{colorlinks=true, urlcolor=primary}

\begin{document}
\noindent \textbf{Lo\"ic Aron MBASSI EWOLO} \\
ENSEA -- Double Dipl\^ome Ing\'enieur \\
\href{mailto:loicaron.mbassiewolo@ensea.fr}{loicaron.mbassiewolo@ensea.fr} \quad (+33) 7 59 52 33 98 \\
\href{https://github.com/Nameless0l}{github.com/Nameless0l} \quad \href{https://mbassiloic.tech}{mbassiloic.tech}

\vspace{1em}\noindent Cannes, le \today
\vspace{1em}\noindent \textbf{Objet : Candidature Alternance -- Ing\'enieur Développement Logiciel Embarqué Satellites}
\vspace{1.5em}\noindent Madame, Monsieur,

\vspace{0.8em}\noindent J'ai impl\'ement\'e des logiciels embarqu\'es complets en VHDL (FPGA Cyclone V, soft-processeur Nios V, contrôleur I2C, debug JTAG avanc\'e) et d\'evelopp\'e des applications C/C++ temps r\'eel sur Nvidia Jetson pour traitement de données Lidar 3D dans le Challenge CoHoMa (ENSEA/Armée française). Formation double diplôme ENSEA/Polytechnique Yaoundé, spécialité systèmes embarqués critiques et électronique temps réel. Stage INRIA Saclay en cours sur l'analyse statique de systèmes logiciels embarqués.

\vspace{0.8em}\noindent Thales Alenia Space, co-constructeur satellites Sentinel et leader européen du spatial, offre l'opportunité unique de développer des logiciels embarqués critiques à bord des satellites pour des programmes majeurs comme Galileo et OneWeb. Mon expérience en architecture logicielle robuste, optimisation temps réel et validation rigoureuse des systèmes critiques correspond directement aux défis du on-board software engineering spatial.

\vspace{0.8em}\noindent J'ai acquis une maîtrise approfondie des protocoles temps réel (I2C, SPI, UART) et des systèmes embarqués critiques à travers le projet Harmony Gloves (capteurs inertiels MPU6050, STM32, protocoles bas niveau) et la modélisation dynamique complète d'un drone avec détection de défauts (Fault Tolerant Control) en MATLAB/Simulink. Ma contribution au projet Jetson CoHoMa démontre ma capacité à optimiser des logiciels C/C++ pour architectures ARM contraintes en ressources et temps réel.

\vspace{0.8em}\noindent Je suis disponible en alternance dès septembre 2026 pour une durée de 12 mois, mobile en France (région Cannes). Je serais ravi de discuter de comment mes compétences en logiciel embarqué C/C++, FPGA/RTL et systèmes temps réel critiques peuvent contribuer aux innovations spatiales de Thales Alenia Space pour les générations futures de satellites européens.

\vspace{1.5em}\noindent Veuillez agr\'eer, Madame, Monsieur, l'expression de mes salutations distingu\'ees.
\vspace{2em}\noindent \textbf{Lo\"ic Aron MBASSI EWOLO}
\end{document}
